% Options for packages loaded elsewhere
\PassOptionsToPackage{unicode}{hyperref}
\PassOptionsToPackage{hyphens}{url}
\documentclass[
]{article}
\usepackage{xcolor}
\usepackage[margin=1in]{geometry}
\usepackage{amsmath,amssymb}
\setcounter{secnumdepth}{5}
\usepackage{iftex}
\ifPDFTeX
  \usepackage[T1]{fontenc}
  \usepackage[utf8]{inputenc}
  \usepackage{textcomp} % provide euro and other symbols
\else % if luatex or xetex
  \usepackage{unicode-math} % this also loads fontspec
  \defaultfontfeatures{Scale=MatchLowercase}
  \defaultfontfeatures[\rmfamily]{Ligatures=TeX,Scale=1}
\fi
\usepackage{lmodern}
\ifPDFTeX\else
  % xetex/luatex font selection
\fi
% Use upquote if available, for straight quotes in verbatim environments
\IfFileExists{upquote.sty}{\usepackage{upquote}}{}
\IfFileExists{microtype.sty}{% use microtype if available
  \usepackage[]{microtype}
  \UseMicrotypeSet[protrusion]{basicmath} % disable protrusion for tt fonts
}{}
\makeatletter
\@ifundefined{KOMAClassName}{% if non-KOMA class
  \IfFileExists{parskip.sty}{%
    \usepackage{parskip}
  }{% else
    \setlength{\parindent}{0pt}
    \setlength{\parskip}{6pt plus 2pt minus 1pt}}
}{% if KOMA class
  \KOMAoptions{parskip=half}}
\makeatother
\usepackage{longtable,booktabs,array}
\newcounter{none} % for unnumbered tables
\usepackage{calc} % for calculating minipage widths
% Correct order of tables after \paragraph or \subparagraph
\usepackage{etoolbox}
\makeatletter
\patchcmd\longtable{\par}{\if@noskipsec\mbox{}\fi\par}{}{}
\makeatother
% Allow footnotes in longtable head/foot
\IfFileExists{footnotehyper.sty}{\usepackage{footnotehyper}}{\usepackage{footnote}}
\makesavenoteenv{longtable}
\usepackage{graphicx}
\makeatletter
\newsavebox\pandoc@box
\newcommand*\pandocbounded[1]{% scales image to fit in text height/width
  \sbox\pandoc@box{#1}%
  \Gscale@div\@tempa{\textheight}{\dimexpr\ht\pandoc@box+\dp\pandoc@box\relax}%
  \Gscale@div\@tempb{\linewidth}{\wd\pandoc@box}%
  \ifdim\@tempb\p@<\@tempa\p@\let\@tempa\@tempb\fi% select the smaller of both
  \ifdim\@tempa\p@<\p@\scalebox{\@tempa}{\usebox\pandoc@box}%
  \else\usebox{\pandoc@box}%
  \fi%
}
% Set default figure placement to htbp
\def\fps@figure{htbp}
\makeatother
\setlength{\emergencystretch}{3em} % prevent overfull lines
\providecommand{\tightlist}{%
  \setlength{\itemsep}{0pt}\setlength{\parskip}{0pt}}
\usepackage{titling}
\pretitle{\begin{center}\Large\bfseries}
\posttitle{\end{center}}
\usepackage{booktabs}
\usepackage{float}
\usepackage{placeins}
\usepackage{graphicx}
\usepackage{amsmath}
\usepackage{amssymb}
\usepackage[font=small]{caption}
\usepackage{bookmark}
\IfFileExists{xurl.sty}{\usepackage{xurl}}{} % add URL line breaks if available
\urlstyle{same}
\hypersetup{
  pdftitle={Posterior-Based Decision Stability under Cross-City Transfer in Urban Flood Risk Prioritisation},
  hidelinks,
  pdfcreator={LaTeX via pandoc}}

\title{Posterior-Based Decision Stability under Cross-City Transfer in
Urban Flood Risk Prioritisation}
\author{Mikel Martinez Mugica\\
Independent Researcher\\
https://habnetic.org}
\date{June 2026}

\begin{document}
\maketitle

{
\setcounter{tocdepth}{3}
\tableofcontents
}
\section{Abstract}\label{abstract}

Urban risk prioritisation workflows frequently produce deterministic
rankings despite substantial uncertainty in hazard representation,
exposure estimation, and model specification. While uncertainty is often
analysed within individual modelling components, downstream
prioritisation decisions are less commonly evaluated explicitly as
probabilistic objects.

This paper presents a Bayesian framework for analysing decision
stability under uncertainty propagation. Rather than treating rankings
as fixed outputs, the framework derives posterior-based decision
quantities including top-k membership probability and borderline
decision share.

A reproducible exposure–hazard–outcome pipeline is implemented using a
logistic Bayesian baseline with simplified exposure and pluvial hazard
proxies. Posterior inference is performed using PyMC and prioritisation
decisions are derived directly from posterior draws. The objective is
methodological evaluation of posterior-derived decision behaviour rather
than hydraulic calibration.

The framework is evaluated across three phases: a Rotterdam baseline,
controlled hazard perturbation experiments, and fixed-specification
cross-city transfer to Hamburg and Donostia–San Sebastián.

Results indicate that decision instability is not diffuse across the
system. Instead, instability remains concentrated near a narrow decision
boundary while most assets exhibit stable prioritisation behaviour
across posterior draws. At comparable 1\% prioritisation thresholds, the
unstable boundary represents 0.0136\% of assets in Rotterdam, 0.0164\%
in Hamburg, and 0.1676\% in Donostia–San Sebastián.

These results demonstrate how posterior inference can be extended beyond
predictive estimation toward explicit analysis of posterior decision
stability under uncertainty.

\begin{center}\rule{0.5\linewidth}{0.5pt}\end{center}

\section{Introduction}\label{introduction}

Urban resilience planning frequently requires prioritising assets for
intervention, such as buildings, infrastructure segments, or
neighbourhoods. These priorities are commonly derived from deterministic
risk scores or expected damage estimates. In practice, they directly
determine which assets receive intervention under limited budgets.
Decision stability is therefore operationally important, but it is often
less directly examined than model-level uncertainty.

Risk estimation involves substantial uncertainty arising from:

\begin{itemize}
\tightlist
\item
  hazard variability
\item
  exposure measurement error
\item
  model specification
\item
  limited observational data
\end{itemize}

Although uncertainty may be analysed within individual modelling
components, it is less commonly propagated through the full modelling
chain to the final prioritisation decisions.

As a result, the stability of prioritisation decisions remains
comparatively underexplored relative to uncertainty analysis within
individual modelling components. This paper proposes treating posterior
decision stability itself as the primary inferential object.

Conventional risk prioritisation asks:

\begin{quote}
Which assets appear most at risk?
\end{quote}

This paper instead asks:

\begin{quote}
Which prioritisation decisions remain reliable once uncertainty has been
propagated through the modelling pipeline?
\end{quote}

Under posterior uncertainty, rankings induce probabilistic membership in
prioritised subsets rather than fixed deterministic decisions. The
present work therefore evaluates prioritisation decisions as
posterior-derived quantities rather than deterministic outputs.

The framework is evaluated across three phases: a Rotterdam baseline,
hazard perturbation experiments, and fixed-specification transfer to
Hamburg and Donostia–San Sebastián. The objective is not predictive
optimisation for individual cities, but to evaluate whether
posterior-derived decision stability remains structurally concentrated
under uncertainty and distributional shift.

\subsection{Contributions}\label{contributions}

This paper makes four main contributions:

\begin{itemize}
\tightlist
\item
  A reproducible Bayesian framework for posterior decision-stability
  analysis
\item
  A set of decision-oriented posterior metrics including top-k
  membership probability and borderline decision share
\item
  A perturbation analysis showing that instability remains localised
  under perturbed hazard representation
\item
  A fixed-specification cross-city transfer experiment for evaluating
  decision stability under distributional shift
\end{itemize}

\begin{center}\rule{0.5\linewidth}{0.5pt}\end{center}

\section{Conceptual Framework}\label{conceptual-framework}

Urban flood risk assessments typically follow a conceptual structure
linking:

Exposure → Hazard → Impact (Outcome)

where exposure describes asset characteristics, hazard represents
environmental forcing, and impact represents realised damage or
disruption.

Within the proposed framework, this structure is extended into a
Bayesian inference workflow:

Exposure → Hazard → Impact → Posterior inference → Posterior decision
metrics

Posterior uncertainty over model parameters induces uncertainty over
asset rankings, which in turn induces uncertainty over prioritisation
membership. Rather than collapsing posterior distributions into
deterministic rankings, the proposed framework derives posterior
decision metrics directly from posterior draws.

\begin{figure}
\centering
\pandocbounded{\includegraphics[keepaspectratio,alt={Hierarchical exposure–hazard–impact model}]{./media-8069bccc9228b131/figures/fig01_graphical_model.pdf}}
\caption{Hierarchical exposure–hazard–impact model}
\end{figure}

Figure 1 summarises the probabilistic structure of the proposed
framework. Exposure indicators and hazard processes jointly determine
impact, posterior inference propagates uncertainty through the model,
and posterior-derived decision metrics quantify uncertainty in
prioritisation decisions. The empirical implementation uses simplified
exposure and pluvial hazard proxies together with a synthetic binary
outcome to isolate decision behaviour under uncertainty. The figure is
intended as a conceptual representation of the proposed framework rather
than a hydraulic flood-process diagram.

\begin{center}\rule{0.5\linewidth}{0.5pt}\end{center}

\section{Experimental Design}\label{experimental-design}

\subsection{Phase 1 — Rotterdam
baseline}\label{phase-1-rotterdam-baseline}

The baseline implementation is performed using Rotterdam as the
reference study area, with 221,324 buildings.

Exposure is represented using hydrographic proximity indicators,
including distance to water and local water-density aggregation. Hazard
is represented using an ERA5-derived pluvial precipitation proxy.

A synthetic Bernoulli outcome variable is used to isolate the behaviour
of the probabilistic inference-to-decision pipeline from empirical
calibration challenges. This means the experiment evaluates the
stability of posterior-derived prioritisation under a controlled outcome
construction; it does not claim validation against observed flood
damage.

\subsection{Phase 2 — Hazard
perturbation}\label{phase-2-hazard-perturbation}

Robustness is evaluated through controlled perturbation of the hazard
proxy.

The standardised hazard representation is perturbed with additive
Gaussian noise while maintaining the same model structure and inference
procedure.

The objective is not to simulate physically realistic perturbations, but
to evaluate whether instability diffuses across the prioritisation
system under perturbed hazard representation.

\subsection{Phase 3 — Fixed-specification cross-city
transfer}\label{phase-3-fixed-specification-cross-city-transfer}

The framework is transferred from Rotterdam to Hamburg and Donostia–San
Sebastián under fixed specification:

\begin{itemize}
\tightlist
\item
  same priors
\item
  same model structure
\item
  same feature definitions
\item
  same scaling reference
\item
  same decision metrics
\end{itemize}

Feature standardisation is performed using Rotterdam statistics, which
are retained as the scaling reference for the transferred cities.

The objective is not predictive optimisation for each city individually,
but evaluation of whether the decision-stability structure persists
under fixed specification and distributional shift.

\begin{center}\rule{0.5\linewidth}{0.5pt}\end{center}

\section{Bayesian Model}\label{bayesian-model}

The model defines a simple generative structure linking exposure, hazard
proxy, and binary impact outcome.

Asset-level impact probability is modelled using a logistic regression
structure:

\[
p_i =
\operatorname{logit}^{-1}
\left(
\alpha + \beta_E E_i + \beta_H H_i
\right)
\]

\[
Y_i \sim \operatorname{Bernoulli}(p_i)
\]

where:

\begin{itemize}
\tightlist
\item
  \(E_i\) denotes the exposure proxy
\item
  \(H_i\) denotes the hazard proxy
\item
  \(Y_i \in \{0,1\}\) denotes the binary impact indicator
\end{itemize}

\subsection{Priors}\label{priors}

Weakly informative Gaussian priors are assigned to the intercept and
regression coefficients:

\[
\alpha, \beta_E, \beta_H \sim \mathcal{N}(0, 2.5)
\]

where 2.5 denotes the prior standard deviation.

\subsection{Inference}\label{inference}

Posterior inference is performed using PyMC with Hamiltonian Monte Carlo
sampling through the NUTS sampler.

Inference is performed using:

\begin{itemize}
\tightlist
\item
  4 chains
\item
  1000 tuning iterations
\item
  1000 posterior draws per chain
\item
  \texttt{target\_accept\ =\ 0.9}
\end{itemize}

Diagnostics include:

\begin{itemize}
\tightlist
\item
  \(\hat{R}\) convergence diagnostics
\item
  effective sample size (ESS)
\item
  prior predictive checks
\item
  posterior predictive checks
\end{itemize}

Prior predictive checks were used to verify that the prior specification
does not imply implausibly high citywide event rates before observing
the synthetic outcome data. Posterior predictive checks were used as
internal model checks under the synthetic Bernoulli outcome
construction. These diagnostics support stable inference behaviour under
the current specification, but they should not be interpreted as
empirical flood validation.

\begin{figure}
\centering
\pandocbounded{\includegraphics[keepaspectratio,alt={Prior predictive event-rate check}]{./media-8069bccc9228b131/figures/fig_prior_predictive_event_rate.pdf}}
\caption{Prior predictive event-rate check}
\end{figure}

Figure 2 shows the prior predictive event-rate check across Rotterdam,
Hamburg, and Donostia–San Sebastián. The prior predictive intervals
remain broad, as expected under weakly informative priors, but low
baseline event rates remain plausible before observing the synthetic
outcome data.

\subsection{Model scope}\label{model-scope}

The logistic specification is intentionally minimal. It serves as a
baseline to isolate the behaviour of posterior-derived decision metrics
rather than to optimise predictive performance.

The present results should therefore not be interpreted as hydraulic
validation or operational flood prediction.

\begin{center}\rule{0.5\linewidth}{0.5pt}\end{center}

\section{Posterior-Derived Decision
Metrics}\label{posterior-derived-decision-metrics}

Rather than relying on point estimates, prioritisation decisions are
derived directly from posterior draws.

For each posterior draw \(s\):

\begin{enumerate}
\def\labelenumi{\arabic{enumi}.}
\tightlist
\item
  asset-level impact probabilities \(p_i^{(s)}\) are computed
\item
  assets are ranked by \(p_i^{(s)}\)
\item
  the top-k subset \(\text{Top}_k^{(s)}\) is extracted
\end{enumerate}

This produces a posterior distribution over prioritisation membership
rather than a single deterministic ranking.

\subsection{Posterior mean risk}\label{posterior-mean-risk}

Posterior mean risk is defined as the posterior expectation of the
asset-level impact probability:

\[
p_{\text{mean}, i}
=
\mathbb{E}_{p(\theta \mid \text{data})}
\left[
P(Y_i = 1 \mid \theta)
\right]
\]

Equivalently, using posterior draws:

\[
p_{\text{mean}, i}
\approx
\frac{1}{S}
\sum_{s=1}^{S}
p_i^{(s)}
\]

\subsection{Top-k membership
probability}\label{top-k-membership-probability}

For a prioritisation threshold \(k\), the posterior probability that
asset \(i\) belongs to the top-k highest-risk assets is:

\[
\pi_{i,k}
=
P(i \in \text{Top}_k \mid \text{data})
\]

This is estimated from posterior draws as:

\[
\pi_{i,k}
\approx
\frac{1}{S}
\sum_{s=1}^{S}
\mathbb{1}
\left(
i \in \text{Top}_k^{(s)}
\right)
\]

This quantity forms the primary inferential object of the analysis.

\subsection{Decision stability
classes}\label{decision-stability-classes}

For a given threshold \(k\), each asset is assigned to one of three
posterior decision-stability classes using \(\pi_{i,k}\):

\[
\pi_{i,k} \leq 0.2
\]

Stable low-priority: the asset is unlikely to belong to the prioritised
set.

\[
0.2 < \pi_{i,k} < 0.8
\]

Unstable boundary: the asset has materially uncertain prioritisation
membership.

\[
\pi_{i,k} \geq 0.8
\]

Stable high-priority: the asset is likely to belong to the prioritised
set.

\subsection{Borderline share}\label{borderline-share}

Borderline share is defined as the proportion of assets in the unstable
boundary class:

\[
\frac{1}{N}\sum_{i=1}^{N}
\mathbb{1}
\left(
0.2 < \pi_{i,k} < 0.8
\right)
\]

This quantity measures how much of the prioritisation system remains
decision-unstable under posterior uncertainty.

\begin{center}\rule{0.5\linewidth}{0.5pt}\end{center}

\section{Results}\label{results}

\subsection{Baseline risk concentration and decision-stability
structure}\label{baseline-risk-concentration-and-decision-stability-structure}

Figure 3 shows the deterministic ranking structure induced by posterior
mean risk in the Rotterdam baseline. Expected risk is concentrated in a
small upper tail rather than distributed evenly across the asset
population.

\begin{figure}[H]
\centering
\includegraphics[width=0.82\textwidth]{figures/phase3_expected_risk_ranking_RTM.pdf}
\caption{Expected-risk ranking for the Rotterdam baseline.}
\end{figure}

This ranking alone is not the main inferential object of the paper. It
motivates the decision-stability analysis by showing where
prioritisation pressure concentrates under a limited intervention
budget.

Figure 4 shows the posterior top-k membership structure for the
Rotterdam baseline at an approximately 1\% prioritisation threshold. The
curve separates the asset population into three decision-stability
regions: stable high-priority assets, an unstable boundary, and stable
low-priority assets.

\begin{figure}[H]
\centering
\includegraphics[width=0.82\textwidth]{figures/phase3_decision_stability_structure_RTM.pdf}
\caption{Posterior decision-stability structure for the Rotterdam baseline.}
\end{figure}

The result is not that all risk estimates are certain. The result is
more specific: uncertainty affecting prioritisation membership is
concentrated near the decision boundary. Most buildings are consistently
classified as either inside or outside the prioritised set across
posterior draws.

\FloatBarrier

\subsection{Decision-stability composition across
cities}\label{decision-stability-composition-across-cities}

Figure 5 compares the share of assets in each decision-stability class
across Rotterdam, Hamburg, and Donostia–San Sebastián at comparable 1\%
prioritisation thresholds.

\begin{figure}[H]
\centering
\includegraphics[width=0.82\textwidth]{figures/phase3_certainty_composition_bar.pdf}
\caption{Decision-stability composition across the three study areas.}
\end{figure}

At comparable thresholds, the unstable boundary remains small in all
three cities. Rotterdam and Hamburg exhibit almost negligible unstable
shares, while Donostia–San Sebastián shows a wider unstable boundary
under fixed-specification transfer. Even in the strongest
transfer-stress case, however, the unstable set remains below 0.2\% of
assets at the comparable 1\% threshold.

\subsection{Robustness to hazard
perturbation}\label{robustness-to-hazard-perturbation}

To evaluate whether the observed concentration of instability depends on
deterministic hazard specification, controlled perturbations are
introduced into the standardised hazard representation:

\[
H_i^{\text{perturbed}} = H_i + \epsilon_i,
\quad
\epsilon_i \sim \mathcal{N}(0, \sigma)
\]

The following perturbation levels are evaluated:

\begin{itemize}
\tightlist
\item
  \(\sigma = 0.00\)
\item
  \(\sigma = 0.05\)
\item
  \(\sigma = 0.10\)
\item
  \(\sigma = 0.20\)
\item
  \(\sigma = 0.30\)
\end{itemize}

{\def\LTcaptype{none} % do not increment counter
\begin{longtable}[]{@{}lr@{}}
\toprule\noalign{}
\(\sigma\) & Borderline share \\
\midrule\noalign{}
\endhead
\bottomrule\noalign{}
\endlastfoot
0.00 & 1.58\% \\
0.05 & 1.30\% \\
0.10 & 1.48\% \\
0.20 & 1.72\% \\
0.30 & 1.74\% \\
\end{longtable}
}

Figure 6 shows that increasing hazard perturbation does not produce
diffuse instability across the system. Instead, instability remains
concentrated near the prioritisation boundary across all tested
perturbation levels.

\begin{figure}[H]
\centering
\includegraphics[width=0.82\textwidth]{figures/fig07_borderline_vs_sigma.pdf}
\caption{Decision instability under hazard perturbation.}
\end{figure}

The perturbation experiment therefore supports the interpretation that
decision instability is structurally localised rather than uniformly
distributed across the asset population.

\subsection{Cross-city transfer}\label{cross-city-transfer}

The same framework is applied to Hamburg and Donostia–San Sebastián
under fixed specification.

To make the transfer comparison interpretable across differently sized
cities, the main cross-city comparison uses approximately 1\%
prioritisation thresholds.

\begin{table}[H]
\centering
\begin{tabular}{lrrrr}
\toprule
City & Representative $k$ & Assets & Prioritised share & Borderline share \\
\midrule
RTM & 2,213 & 221,324 & 1.00\% & 0.0136\% \\
HAM & 3,415 & 341,530 & 1.00\% & 0.0164\% \\
DON & 78 & 7,755 & 1.01\% & 0.1676\% \\
\bottomrule
\end{tabular}
\caption{Cross-city comparison at approximately 1\% prioritisation thresholds.}
\end{table}

Table 1 summarises the main transfer result. Rotterdam and Hamburg
exhibit almost negligible borderline share. Donostia–San Sebastián shows
a wider unstable boundary under fixed-specification transfer, but even
there the unstable boundary remains below 0.2\% of assets at the
comparable 1\% prioritisation threshold.

Figure 7 shows the same comparison using log-scaled rank, which makes
the narrow transition region near the prioritisation boundary easier to
inspect.

\begin{figure}[H]
\centering
\includegraphics[width=0.82\textwidth]{figures/phase3_cross_city_stability_structure.pdf}
\caption{Cross-city posterior top-k membership structure using log-scaled rank.}
\end{figure}

Across all three cities, posterior instability remains concentrated near
a narrow prioritisation boundary. Rotterdam and Hamburg exhibit highly
polarised membership structure, while Donostia–San Sebastián exhibits
moderate local broadening of the unstable boundary region consistent
with distributional shift.

Figure 8 shows how borderline share evolves across prioritisation
thresholds for all three cities.

\begin{figure}[H]
\centering
\includegraphics[width=0.82\textwidth]{figures/phase3_borderline_share_vs_k.pdf}
\caption{Borderline share as a function of the prioritisation threshold.}
\end{figure}

Across all evaluated thresholds, instability remains concentrated in a
relatively small subset of assets. Donostia–San Sebastián exhibits a
wider unstable boundary under fixed-specification transfer, but the
overall instability structure remains highly localised relative to the
full asset population.

\subsection{Spatial structure of posterior decision
stability}\label{spatial-structure-of-posterior-decision-stability}

Figure 9 shows the spatial structure of posterior top-k membership
probability and local transition regions for Rotterdam, Hamburg, and
Donostia–San Sebastián under comparable prioritisation thresholds.

The spatial comparison illustrates that posterior instability is not
spatially diffuse across the urban system. Instead, uncertainty in
prioritisation membership remains concentrated within relatively narrow
local transition structures separating stable high-priority and stable
low-priority assets.

Rotterdam and Hamburg exhibit highly polarised posterior membership
distributions with extremely narrow unstable regions. Donostia–San
Sebastián exhibits moderate local broadening of the transition structure
under fixed-specification transfer, consistent with stronger
distributional shift relative to the Rotterdam reference specification.

These spatial figures are not presented as hydraulic validation. Their
purpose is methodological: they expose the spatial organisation of
posterior-derived decision stability under uncertainty propagation and
cross-city transfer.

\clearpage

\begin{figure}[!p]
\centering
\includegraphics[width=0.95\textwidth]{figures/RTM_paper_citywide_topk_map.pdf}

\vspace{0.5em}

\includegraphics[width=0.82\textwidth]{figures/RTM_paper_boundary_zoom_map.pdf}

\caption{Spatial structure of posterior top-k membership probability and local transition region for Rotterdam. The citywide map shows posterior top-k membership probability; the zoom map shows the local transition region near the prioritisation boundary.}
\end{figure}

\clearpage

\begin{figure}[!p]
\centering
\includegraphics[width=0.95\textwidth]{figures/HAM_paper_citywide_topk_map.pdf}

\vspace{0.5em}

\includegraphics[width=0.82\textwidth]{figures/HAM_paper_boundary_zoom_map.pdf}

\caption{Spatial structure of posterior top-k membership probability and local transition region for Hamburg. The citywide map shows posterior top-k membership probability; the zoom map shows the local transition region near the prioritisation boundary.}
\end{figure}

\clearpage

\begin{figure}[!p]
\centering
\includegraphics[width=0.95\textwidth]{figures/DON_paper_citywide_topk_map.pdf}

\vspace{0.5em}

\includegraphics[width=0.82\textwidth]{figures/DON_paper_boundary_zoom_map.pdf}

\caption{Spatial structure of posterior top-k membership probability and local transition region for Donostia--San Sebastián. The citywide map shows posterior top-k membership probability; the zoom map shows the local transition region near the prioritisation boundary.}
\end{figure}

\clearpage

\begin{center}\rule{0.5\linewidth}{0.5pt}\end{center}

\section{Discussion}\label{discussion}

This study demonstrates how prioritisation decisions can be analysed as
probabilistic objects rather than deterministic rankings. By deriving
decision metrics directly from posterior distributions, the proposed
framework evaluates not only which assets appear most at risk, but which
prioritisation decisions remain reliable under uncertainty.

Across the baseline, perturbation, and transfer experiments, posterior
uncertainty remains concentrated near a relatively narrow decision
boundary while most assets exhibit stable prioritisation behaviour
across posterior draws. These findings suggest that uncertainty
propagation does not necessarily imply diffuse decision instability.
Instead, uncertainty can remain localised near prioritisation
thresholds, allowing most decisions to remain stable even when
predictive uncertainty is substantial.

The perturbation experiments further suggest that this concentration of
instability is not highly sensitive to moderate changes in the hazard
representation. Similarly, the fixed-specification transfer experiments
indicate that the overall decision-stability structure may persist under
moderate distributional shift, although the width of the unstable
boundary varies across cities.

These findings should not be interpreted as evidence of predictive
generalisation or hydraulic validity across study areas. Rather, the
experiments evaluate the structural behaviour of posterior-derived
decision metrics under a fixed probabilistic specification. The
framework therefore assesses the stability of prioritisation decisions
conditional on the assumed model, rather than the empirical correctness
of the prioritisation itself.

\begin{center}\rule{0.5\linewidth}{0.5pt}\end{center}

\section{Limitations}\label{limitations}

Several limitations should be noted.

\begin{itemize}
\tightlist
\item
  The hazard representation is proxy-based and does not model hydraulic
  flood dynamics.
\item
  The outcome variable is synthetic and not calibrated against observed
  damage data.
\item
  Exposure is represented using simplified hydrographic proximity
  indicators.
\item
  The current framework evaluates a single model family with limited
  prior sensitivity analysis.
\item
  The framework does not currently include latent hazard processes or
  utility-theoretic decision modelling.
\end{itemize}

These limitations are central to the interpretation of the results. The
analysis does not claim operational flood prediction, hydraulic
validation, or empirical correctness of the prioritised assets. It
isolates the statistical behaviour of posterior-derived prioritisation
metrics under a deliberately simplified model.

No claims of operational deployment or empirical flood prediction are
made.

\begin{center}\rule{0.5\linewidth}{0.5pt}\end{center}

\section{Conclusion}\label{conclusion}

This paper presented a Bayesian framework for analysing posterior
decision stability under uncertainty propagation.

Rather than treating rankings as deterministic outputs, the framework
derives posterior-based decision quantities that quantify uncertainty in
prioritisation membership itself.

Across baseline, perturbation, and transfer experiments, instability
remains concentrated near a narrow prioritisation boundary while most
assets exhibit stable prioritisation behaviour under posterior
uncertainty. Under comparable 1\% prioritisation thresholds, unstable
boundary shares remain small across all three evaluated cities, even
under fixed-specification transfer.

The contribution is methodological rather than hydraulic. The results
demonstrate how posterior inference can be extended from predictive
estimation toward explicit analysis of posterior decision stability
under uncertainty.

\begin{center}\rule{0.5\linewidth}{0.5pt}\end{center}

\section{References}\label{references}

Hall, J. W., \& Harvey, H. (2009). Decision making under severe
uncertainties for flood risk management: A case study of Info-Gap
robustness analysis. \emph{Journal of Flood Risk Management}.

Lv, H., Wu, Z., Guan, X., \& Meng, Y. (2021). The construction of flood
loss ratio function in cities lacking loss data based on dynamic
proportional substitution and hierarchical Bayesian model. \emph{Journal
of Hydrology}.

McMillan, H., \& Brasington, J. (2008). End-to-end flood risk
assessment: A coupled model cascade with uncertainty estimation.
\emph{Water Resources Research}.

Mohor, G. S., et al.~(2021). Residential flood loss estimated from
Bayesian multilevel models. \emph{Natural Hazards and Earth System
Sciences}.

Sairam, N., Schröter, K., Rözer, V., Merz, B., \& Kreibich, H. (2019). A
Bayesian hierarchical model for flood damage estimation in data-scarce
regions. \emph{Water Resources Research}.

Wu, Y., et al.~(2019). Assessing urban flood disaster risk using
Bayesian network model and GIS applications. \emph{International Journal
of River Basin Management}.

Wu, Y., et al.~(2020). Urban flood disaster risk evaluation based on
ontology and Bayesian Network. \emph{Journal of Hydrology}.

\section{Appendix A. Model input
sample}\label{appendix-a.-model-input-sample}

A five-row sample of the Rotterdam model input table is provided to make
the asset-level schema explicit.

The complete processed input tables are stored in:

\begin{itemize}
\tightlist
\item
  \texttt{outputs/phase3/RTM/phase3\_features\_scaled.parquet}
\item
  \texttt{outputs/phase3/HAM/phase3\_features\_scaled.parquet}
\item
  \texttt{outputs/phase3/DON/phase3\_features\_scaled.parquet}
\end{itemize}

\begin{table}[H]
\centering
\small
\begin{tabular}{rrrrr}
\toprule
Building ID & Exposure ($E$) & Hazard mm & Hazard log-rel & damage \\
\midrule
305012 & -0.033 & 25.422 & -0.00868 & 0 \\
313960 & 0.238 & 25.419 & -0.00881 & 0 \\
313263 & -0.131 & 25.423 & -0.00864 & 0 \\
310491 & -0.273 & 25.425 & -0.00859 & 0 \\
313127 & -0.342 & 25.424 & -0.00862 & 0 \\
\bottomrule
\end{tabular}
\end{table}

The sample is included for transparency only and is not used directly
for inference beyond illustrating the asset-level schema consumed by the
model.

\end{document}
